\chapter{Discussion and Conclusions}
\label{ch:discussion_conclusions}

\section{Chapter Overview}
\label{sec:concl_overview}
In this chapter, a summary of the main results, practical achievements, lessons learned, and directions for future research concerning the proposed paper-based customizable virtual keyboard is provided. By relying only on a single webcam and plain paper without specialized devices like depth cameras, Time-of-Flight sensors, laser projectors, or sensor gloves, this project demonstrates that a normal sheet of paper can serve as an interactive and customizable virtual keyboard running in real time on an ordinary computer CPU.

Section~\ref{sec:concl_objectives} discusses how each of the five specific research objectives was accomplished by triangulating goals, system implementation, and experimental test results. Section~\ref{sec:concl_problems} explains the practical problems faced during the project and how they were solved. Section~\ref{sec:concl_reflection} offers personal reflections on the learning process and skills gained. Section~\ref{sec:concl_business} discusses practical real-world application possibilities and benefits. Finally, Section~\ref{sec:concl_future} outlines future research directions, and Section~\ref{sec:concl_summary} concludes the chapter.

\section{Accomplishment of Research Objectives}
\label{sec:concl_objectives}
To evaluate the research outcomes objectively, we use the method of triangulation. Triangulation links the experimental results obtained in Chapter~\ref{ch:results} with the original research objectives stated in Chapter~\ref{chap:objectives}, supported by the system implementation and algorithms described in Chapter~\ref{ch:methodology}.

\subsection{Triangulation Evaluation of Objective 1: Theoretical Analysis and Baseline Requirements}
\label{subsec:concl_obj1}
\begin{itemize}
    \item \textbf{Objective Statement:} \emph{To review and analyze existing visual input systems, fiducial marker tracking methods, planar homography algorithms, and hand pose estimation models, establishing baseline technical requirements for monocular touch detection on flat printed paper.}
    \item \textbf{Triangulation Evidence:} Chapter~\ref{ch:literature_review} reviewed prior virtual keyboard literature across eight sensing modalities and summarized them in a structured taxonomy (Figure~\ref{fig:lit_conceptual_map_diagram}) and comparative table (Table~\ref{tab:lit_comparative_matrix}). Six key research gaps were identified: (1) multi-finger tracking restrictions in single cameras, (2) artificial dwell-time delays, (3) heavy GPU computing requirements, (4) brittleness under ambient lighting and shadows, (5) dependence on expensive depth hardware, and (6) rigid hard-coded layouts.
    \item \textbf{Validation Outcome:} \textbf{Objective Fully Accomplished.} The literature taxonomy directly established the technical requirements for the software architecture and engineering pipeline.
\end{itemize}

\subsection{Triangulation Evaluation of Objective 2: Hand Kinematic Representation and Scale Normalization}
\label{subsec:concl_obj2}
\begin{itemize}
    \item \textbf{Objective Statement:} \emph{To investigate the kinematic properties of hand joints during typing, formulate unitless hand-length scale normalization ($L_{\text{hand}}$) invariant to camera distance and resolution, and establish a 12~FPS sliding window pipeline that preserves impact deceleration without temporal smoothing delays.}
    \item \textbf{Triangulation Evidence:} The kinematic data pipeline was designed in Chapter~\ref{ch:methodology} (Section~\ref{sec:meth_data_pipeline}). Landmark coordinates were scale-normalized using hand length $L_{\text{hand}} = \|\mathbf{P}_9 - \mathbf{P}_0\|_2$, and temporal features were assembled into 5-frame sliding windows with 2-frame overlap at 12~FPS. Crucially, direct feature propagation without low-pass filters preserved physical impact decelerations. In Chapter~\ref{ch:results} (Section~\ref{sec:test_results_review}), distance ablation tests confirmed steady accuracy across $30\text{ cm}$ to $100\text{ cm}$ ($F1 \ge 0.958$), and resolution invariance tests showed consistent performance across 360p, 480p, 720p, and 1080p ($F1 \in [0.961, 0.964]$).
    \item \textbf{Validation Outcome:} \textbf{Objective Fully Accomplished.} Scale normalization and direct feature propagation effectively eliminated scale and resolution variations.
\end{itemize}

\subsection{Triangulation Evaluation of Objective 3: Deep Learning Sequence Model Architecture Benchmark}
\label{subsec:concl_obj3}
\begin{itemize}
    \item \textbf{Objective Statement:} \emph{To benchmark, train, and optimize lightweight deep learning sequence models for independent multi-finger touch contact classification, selecting an optimal model that operates under 3~ms on a standard CPU while achieving high touch detection precision and recall across all five digits.}
    \item \textbf{Triangulation Evidence:} A multi-model benchmark suite (\texttt{run\_all.py}) evaluated 22 pure 2D model architecture combinations across five model families (1D CNN, Attention, ResNet, BiLSTM, LSTM). As documented in Chapter~\ref{ch:results} (Table~\ref{tab:model_architecture_comparison}), the uni-directional PyTorch LSTM model trained on normalized position and velocity features achieved an F1-score of $0.963$, an average per-digit F1-score of $0.965$ (Table~\ref{tab:per_digit_confusion_matrix}), and a fast forward pass time of $2.08\text{ ms}$ on CPU with a compact model size of $0.44\text{ MB}$ (\texttt{best\_finger\_touch\_lstm.pth}).
    \item \textbf{Validation Outcome:} \textbf{Objective Fully Accomplished.} The selected LSTM model satisfied all accuracy and CPU speed requirements.
\end{itemize}

\subsection{Triangulation Evaluation of Objective 4: System Architecture Engineering and Two-Application Suite Design}
\label{subsec:concl_obj4}
\begin{itemize}
    \item \textbf{Objective Statement:} \emph{To architect and implement an integrated two-application software suite enabling custom layout design and decoupled action profile multiplexing on the same printed paper sheet.}
    \item \textbf{Triangulation Evidence:} The two-application software suite was implemented as detailed in Chapter~\ref{ch:methodology}: Application 1 (Layout Designer, \texttt{designer\_app.py}) provides a PySide6 drag-and-drop canvas, automated AprilTag border anchors, XML layout exports, and 300 DPI vector PDF printing. Application 2 (Runtime Virtual Keyboard Engine) combines a setup mapper GUI (Component A) for assigning keys and loading JSON action profiles with a background multi-threaded vision engine (Component B) that tracks AprilTags, extracts MediaPipe landmarks, runs PyTorch LSTM inference, and transforms coordinates ($P_{\text{XML}} = H \cdot P_{\text{pixel}}$). Functional tests in Chapter~\ref{ch:results} (Table~\ref{tab:functional_test_cases}) yielded a $100\%$ pass rate across all test cases (FT-01 to FT-10), confirming successful action multiplexing on identical paper sheets.
    \item \textbf{Validation Outcome:} \textbf{Objective Fully Accomplished.} Decoupled layout design and runtime action multiplexing were fully realized.
\end{itemize}

\subsection{Triangulation Evaluation of Objective 5: Empirical System Evaluation and Usability Benchmarking}
\label{subsec:concl_obj5}
\begin{itemize}
    \item \textbf{Objective Statement:} \emph{To experimentally evaluate the complete virtual keyboard pipeline across touch detection accuracy, homography stability under tilt, end-to-end CPU execution latency, typing performance (WPM and CER), and user interaction feedback.}
    \item \textbf{Triangulation Evidence:} Empirical testing in Chapter~\ref{ch:results} confirmed an end-to-end CPU latency of $29.09\text{ ms}$ ($34.4\text{ FPS}$), operating well within the $83.33\text{ ms}$ budget at 12~FPS, with low resource consumption ($142.5\text{ MB}$ RAM, $24.8\%$ single-core CPU load). AprilTag homography tracking maintained an RMS error under $0.42\text{ mm}$ for camera tilts up to $60^\circ$. User typing tests with ten participants achieved an average typing speed of $24.6\text{ WPM}$ and a Character Error Rate of $5.2\%$, substantially outperforming prior single-camera virtual keyboard baselines ($12.8\text{ WPM}$, $11.4\%$ CER \cite{Ji2018Local}).
    \item \textbf{Validation Outcome:} \textbf{Objective Fully Accomplished.} Rigorous benchmarks validated real-time response, tracking stability, and usability.
\end{itemize}

Table~\ref{tab:objective_triangulation_summary} presents the objective accomplishment and triangulation matrix across all five specific objectives.

\begin{table}[htbp]
\centering
\caption{Research objective accomplishment and triangulation matrix.}
\label{tab:objective_triangulation_summary}
\resizebox{\textwidth}{!}{%
\begin{tabular}{p{3.5cm} p{3.5cm} p{4.0cm} p{3.5cm} c}
\toprule
\textbf{Research Objective} & \textbf{Design / Chapter Origin} & \textbf{Implementation Artifact} & \textbf{Empirical Validation Evidence} & \textbf{Status} \\
\midrule
\textbf{O1: Baseline Requirements and Literature Gaps} & Chapter 3 Literature Review & Research DB (\texttt{references.bib}, \texttt{summary-matrix.md}) & Identified 6 key gaps; established evaluation baseline. & \textbf{ACCOMPLISHED} \\
\midrule
\textbf{O2: Scale Normalization and Kinematic Pipeline} & Chapter 4 Methodology & Partition 2 (\texttt{datacreator/}) scripts & Distance invariance ($30$--$100\text{ cm}$); resolution invariance ($360$p--$1080$p). & \textbf{ACCOMPLISHED} \\
\midrule
\textbf{O3: Deep Learning Sequence Architecture Benchmark} & Chapter 4 Methodology & Partition 3 (\texttt{run\_all.py}) benchmark & F1-score $0.963$; CPU inference $2.08\text{ ms}$; model size $0.44\text{ MB}$. & \textbf{ACCOMPLISHED} \\
\midrule
\textbf{O4: Two-Application System Suite Creation} & Chapter 4 Methodology & Application 1 (\texttt{designer\_app.py}) and Application 2 Suite & $100\%$ pass rate across FT-01--FT-10; same-layout action profile switching. & \textbf{ACCOMPLISHED} \\
\midrule
\textbf{O5: Empirical Latency and Typing Usability Evaluation} & Chapter 5 Results & Multi-threaded Engine and Evaluation Scripts & Latency $29.09\text{ ms}$ ($34.4\text{ FPS}$); $24.6\text{ WPM}$; $5.2\%\text{ CER}$; tilt up to $60^\circ$. & \textbf{ACCOMPLISHED} \\
\bottomrule
\end{tabular}%
}
\end{table}

\section{Problems Encountered and Technical Mitigations}
\label{sec:concl_problems}
A number of practical issues arose during the development and testing of the system, requiring specific technical mitigations.

\subsection{Motion Blur from Standard Rolling Shutter Webcams}
Standard commodity webcams capture video at $30\text{ FPS}$ using rolling shutter sensors that produce motion blur during rapid finger taps \cite{Thomas2013Camera}. Quick downward finger motions occasionally caused slight landmark tracking jitter.
\begin{itemize}
    \item \textbf{Technical Mitigation:} Resampling training video feeds to a uniform 12~FPS rate (\texttt{resample\_12fps.py}) trained the PyTorch LSTM network on motion patterns across consistent time intervals. The temporal windowing allows the model to handle slight frame blur during live camera typing without missing key events.
\end{itemize}

\subsection{Lighting Changes, Shadows, and Surface Reflections}
Sudden changes in room lighting, window shadows, and glare on paper surfaces can disrupt image thresholding and marker detection \cite{ReviewVirtualKeyboard2020}. Shadows cast by user hands across the paper can interfere with corner tracking.
\begin{itemize}
    \item \textbf{Technical Mitigation:} The AprilTag detector applies adaptive local thresholding before finding marker boundaries (\texttt{aprilTag/estimater.py}). In addition, placing the four fiducial tags at the outer corners of the sheet ensures that when hands type in the center, at least one pair of markers remains clearly visible. If a marker is momentarily blocked, the system keeps the previous valid homography matrix ($H$) in memory until all markers are visible again.
\end{itemize}

\subsection{Tendon Coupling Between Adjacent Fingers}
Human hand anatomy features shared flexor tendons, especially between the ring and pinky fingers. When a user presses a key firmly with the ring finger, the pinky finger often moves downward involuntarily, which could lead to accidental key presses in simple thresholding systems.
\begin{itemize}
    \item \textbf{Technical Mitigation:} The PyTorch LSTM network evaluates all five fingers simultaneously using a multi-label Binary Cross-Entropy loss ($\mathcal{L}_{\text{BCE}}$). By observing joint angles across the full 21-joint skeleton rather than tracking fingertips alone, the neural network learns to distinguish an intentional finger strike from passive sympathetic movements.
\end{itemize}

\subsection{Absence of Physical Switch Feedback}
Typing on a sheet of paper laid flat on a table provides a solid stopping surface, but there is no mechanical button travel or switch click sound. In early tests, users were sometimes unsure whether their keystroke had registered and pressed down harder than necessary.
\begin{itemize}
    \item \textbf{Technical Mitigation:} Application 2 provides immediate visual and audio feedback whenever a touch is confirmed ($p_k > 0.90$). The system plays a short click sound (latency under $5\text{ ms}$) and highlights the touched key in green on the screen overlay, giving users clear feedback that their action was registered.
\end{itemize}

\section{Self-Reflection}
\label{sec:concl_reflection}
Working on this final year research project provided valuable theoretical knowledge and practical hands-on experience in computer vision, deep learning, and software development.

\subsection{Shifting Views on Vision-Based Touch Interfaces}
Before starting this research, I assumed that creating a responsive virtual keyboard would require expensive hardware accessories, such as Time-of-Flight depth sensors, laser projectors, or sensor gloves. However, studying projective homography and hand kinematics showed me that thoughtful software design can often replace specialized hardware. By combining standard printed paper, a basic webcam, and a lightweight LSTM network running on a commodity CPU, an interactive virtual input surface can be created at virtually zero extra cost. This experience showed me how software solutions can make computer interfaces more accessible in resource-limited settings.

\subsection{Technical Skills Gained Across Domains}
During this project, I gained practical skills across four key areas of computer science:
\begin{enumerate}
    \item \textbf{Computer Vision and Projective Geometry:} Learned camera matrix calibration ($K$), lens distortion correction ($\mathbf{D}$), AprilTag detection, and solving planar homography matrices ($H \in \mathbb{R}^{3 \times 3}$) using the Direct Linear Transform (DLT) and Singular Value Decomposition (SVD).
    \item \textbf{Deep Learning for Time-Series Signals:} Gained hands-on experience in designing and training recurrent neural networks with PyTorch, writing an automated benchmark suite (\texttt{run\_all.py}), implementing unitless scale normalization ($L_{\text{hand}}$), and assembling temporal sliding windows.
    \item \textbf{Multi-Threaded Programming in Python and Qt:} Implemented producer-consumer thread architectures in PySide6 (\texttt{QThread}, \texttt{QMutex}, \texttt{QWaitCondition}) to separate camera capture from neural network inference, preventing UI freezes and maintaining real-time throughput ($34.4\text{ FPS}$).
    \item \textbf{Academic Research and Testing:} Developed skills in designing structured test cases, measuring HCI typing metrics (Words Per Minute and Character Error Rate), and writing a technical dissertation in LaTeX using IEEE reference formatting.
\end{enumerate}

\subsection{Key Challenges and Methodological Reflections}
\label{subsec:concl_learning_curves}
A major learning challenge in this project was balancing computational complexity with real-time response on an ordinary laptop CPU. Initially, my instinct was to process video at 30 to 60~FPS using large deep learning models. However, testing showed that large models caused CPU bottlenecks and high latency without a dedicated GPU, making the system impractical for regular computers.

Through profiling, I discovered that sub-sampling video to a steady 12~FPS rate captured the necessary impact deceleration information while providing an $83.33\text{ ms}$ processing window. This allowed the entire vision and deep learning pipeline to run in just $29.09\text{ ms}$ on a commodity CPU, leaving plenty of room for background tasks. Unitless hand scale normalization and SVD homography also ensured that touch detection accuracy ($F1 \approx 0.963$) remained consistent even when testing with lower-resolution camera streams, such as 360p or 480p.

\section{Practical Applications and Real-World Impact}
\label{sec:concl_business}
The customizable paper virtual keyboard offers several practical benefits over traditional mechanical keyboards, including low production costs, user-defined layouts, easy sanitization, and flexible action mapping.

\subsection{Hygienic Input in Medical and Laboratory Environments}
In hospital wards, clinics, and cleanrooms, regular mechanical keyboards and touchscreens are difficult to clean thoroughly and can harbor germs and pathogens. In contrast, a paper keyboard with AprilTag anchors can be laminated for wipe-down disinfection or simply discarded and replaced with a fresh printout for pennies per sheet, providing a sterile input solution for healthcare staff.

\subsection{Hazardous Cleanrooms and Industrial Environments}
In chemical laboratories and cleanrooms where electrical equipment must be non-sparking, standard keyboards can pose dust or electrical hazards. A plain paper keyboard sheet is completely passive and carries no electrical currents, making it safe to use on work surfaces in sensitive environments.

\subsection{Accessible Classrooms in Developing Regions}
In schools and computer labs in developing areas, purchasing hardware accessories or replacing damaged keyboards is often a financial challenge. Using Application 1 (Layout Designer), teachers can print custom keyboard layouts, math pads, piano keys, or interactive exercise sheets on standard A4 paper. Connected to a budget PC with a basic webcam, students can type and interact directly on printed paper sheets.

\subsection{Custom Macro Controllers and Accessibility Tools}
Hardware macro pads for digital audio workstations (like Ableton Live), video editing suites, or design software are expensive and have fixed keys. With Application 2, users can switch between different software action profiles (for example, recording, editing, or mixing shortcuts) on the exact same printed sheet of paper without reprinting. Additionally, users with motor impairments can design custom layouts with enlarged keys tailored to their range of movement, switching between typing, navigation, and shortcuts on the same physical sheet.

\section{Future Recommendations}
\label{sec:concl_future}
While the system performs reliably, several interesting directions remain for future research:

\subsection{Lightweight Monocular Neural Depth Estimation}
Future research could explore integrating lightweight neural depth estimation models (such as MiDaS-nano or Depth Anything) directly into the pipeline. Combining 2D skeletal joint positions with coarse depth maps could provide millimeter-level Z-axis information to further refine touch contact classification.

\subsection{Multi-Camera Setup for Occlusion Reduction}
Using two webcams placed at different angles instead of a single camera could help solve the hand occlusion problem. Cross-view 3D triangulation would keep tracking active even when user hands or sleeves temporarily block one camera perspective.

\subsection{Two-Handed (Bimanual) Typing Extensions}
While the current system is optimized for single-handed typing to keep CPU usage low and avoid covering fiducial anchors on A4 paper, the modular architecture can be extended to two hands. Future work can set MediaPipe to track two hands (\texttt{num\_hands=2}) with separate sliding window buffers. Because the PyTorch LSTM forward pass takes only $2.08\text{ ms}$ on CPU, dual-hand inference would take approximately $33\text{ ms}$ total, remaining well within the $83.33\text{ ms}$ budget at 12~FPS.

\subsection{Mid-Air Ultrasonic and Acoustic Haptic Feedback}
Ultrasonic haptic arrays could be investigated to provide tactile feedback without mechanical switches. Directing acoustic pressure waves toward the active fingertip upon key press confirmation would provide a physical sensation of contact in mid-air.

\subsection{Deep Learning Based Fiducial Marker Tracking}
While traditional AprilTag tracking runs efficiently on CPU, future work could test compact convolutional tag detectors like DeepTag \cite{Zhang2022DeepTag}. Direct neural keypoint detection could improve homography stability under severe motion blur or partial marker occlusion.

\subsection{Temporal Transformer Architectures}
Although the uni-directional PyTorch LSTM model achieves high accuracy ($F1 = 0.963$), exploring lightweight temporal Transformer models could help model long-term sequential typing habits and further reduce Character Error Rates during fast typing.

\section{Chapter Summary}
\label{sec:concl_summary}
This chapter concluded the thesis by reviewing the achievement of all five research objectives, discussing technical challenges and solutions, reflecting on personal learning experiences, presenting practical application scenarios, and outlining directions for future work.

The research demonstrated that a customizable paper virtual keyboard powered by a single monocular webcam, unitless hand scale normalization ($L_{\text{hand}}$), AprilTag homography ($H$), and a lightweight PyTorch LSTM network (\texttt{best\_finger\_touch\_lstm.pth}) delivers accurate, real-time typing ($29.09\text{ ms}$ latency, $34.4\text{ FPS}$, $24.6\text{ WPM}$) on standard computer CPUs without a GPU. By replacing specialized hardware with software engineering, this work provides an accessible, hygienic, and flexible human-computer interaction interface.
